\section{Shared qualification for complete dynamic policies}\label{sec:recursive86}
The first-date comparison has a natural dynamic extension, but it must account for the distribution of states induced by the policy being evaluated. A change in an early action changes both the continuation rule and the probability of reaching later observation cells. Holding the reference occupation law fixed would omit the second effect. We develop a transfer that retains both effects and the restriction that one parameter vector generates all the decisions.

We first consider the original expected-cost problem of Section~\ref{sec:core49}. Write $K_t(\cdot\mid x,a)$ for its original transition kernel and let $\bar\pi$ be any feasible reference policy. Its complete continuation cost is $\bar V_t$, including all subsequent actions of $\bar\pi$, with $\bar V_T=g$. Define the signed reference advantage
\begin{equation}\label{eq:advantage86}
 d_t(x,a)=c_t(x,a)+\beta_t K_t\bar V_{t+1}(x,a)-\bar V_t(x).
\end{equation}
No optimal value is needed in this definition. For a qualification class $\Theta_\epsilon$, let $\pi^\theta$ be the complete feasible policy generated by $\theta$, including its observation and tie rules. Its date-$t$ state--action occupation measure is $m_t^\theta$. Both its actions and its occupation measures may differ from the reference at every date.

\begin{theorem}[Dynamic shared-parameter transfer]\label{thm:dynamic86}
Suppose the costs are bounded and the policies and kernels are measurable. For every $\theta\in\Theta_\epsilon$,
\begin{equation}\label{eq:occupation86}
 J_0^{\pi^\theta}-J_0^{\bar\pi}
 =\sum_{t<T}B_t\int d_t(x,a)\,m_t^\theta(dx,da).
\end{equation}
Let $\mathcal R_\epsilon$ be any set of occupation measures that contains those generated by all $\theta\in\Theta_\epsilon$ and satisfies the original-law flow equations. If $d_t^-\le d_t\le d_t^+$, then
\begin{equation}\label{eq:dynamic-bands86}
 \inf_{m\in\mathcal R_\epsilon}\sum_{t<T}B_t\int d_t^-\,dm_t
 \ \le\ J_0^{\pi^\theta}-J_0^{\bar\pi}\ \le\
 \sup_{m\in\mathcal R_\epsilon}\sum_{t<T}B_t\int d_t^+\,dm_t.
\end{equation}
The inequalities hold simultaneously for the entire qualification class. They do not require a common continuation policy after the first date, a common state occupation law, or a finite state space. Valid restrictions induced by the shared parameter vector can be added to $\mathcal R_\epsilon$ without weakening either bound.
\end{theorem}

The identity in \eqref{eq:occupation86} is a policy-performance telescoping identity, rather than a new principle of dynamic programming; related performance comparisons appear in \citet{munos2003,kakade2002}. The substantive qualification problem is to construct an occupation set that is valid for the original law and respects the same fitted parameters at different dates and states. The following finite implementation makes those obligations explicit.

\subsection{An original-law flow relaxation}
At date $t$, an implemented observation rule has cells $C_{t,i}$ and a finite menu $\mathcal A_{t,i}$ of actions feasible throughout the corresponding cell. These are observation cells, not an assumption that the true process is a finite-state Markov chain. Suppose certified constants satisfy, for every $x\in C_{t,i}$,
\[
 d_{t,i,a}^-\le d_t(x,a)\le d_{t,i,a}^+,\qquad
 p_{t,i,a,j}^-\le K_t(C_{t+1,j}\mid x,a)\le p_{t,i,a,j}^+.
\]
The transition bounds must hold throughout the cell. Averaging a kernel under a uniform cell distribution is not sufficient when policy-dependent selection changes that distribution.

Let $x_{t,i,a}$ be an occupation mass, $e_{t,i,a,j}$ a transition mass, and $z_{t,i,a}$ an action-choice indicator. Before relaxation $z$ is binary. Impose
\begin{align}\label{eq:flows86}
 \sum_a x_{0,i,a}&=\mu(C_{0,i}), &
 \sum_a x_{t+1,j,a}&=\sum_{i,a}e_{t,i,a,j},\\
 \sum_j e_{t,i,a,j}&=x_{t,i,a}, &
 p_{t,i,a,j}^-x_{t,i,a}&\le e_{t,i,a,j}\le p_{t,i,a,j}^+x_{t,i,a},\\
 0\le x_{t,i,a}&\le z_{t,i,a}, &
 \sum_a z_{t,i,a}&=1,\quad 0\le z_{t,i,a}\le1.\notag
\end{align}
Transition equations are needed for $t<T-1$; terminal costs are already included in \eqref{eq:advantage86}. Qualification can exclude an action or a joint event involving cells at different dates. If no $\theta\in\Theta_\epsilon$ can select all choices in $E$, add
\begin{equation}\label{eq:flow-conflict86}
 \sum_{(t,i,a)\in E}z_{t,i,a}\le |E|-1.
\end{equation}
A qualification lower certificate strictly exceeding its allowed objective excludes the event. Equality does not. Full-menu inequalities, feasible action repairs and a single coefficient vector are used in generating the event, as in Section~\ref{sec:compatible84}. Unchecked or inconclusive events add no restriction.

Every actual policy defines a feasible binary $z$, its true occupation masses $x$, and its true transition masses $e$. Integrating the cellwise transition inequalities proves \eqref{eq:flows86}, even when the conditional state law within a cell is nonuniform. Replacing binary choices by the displayed unit intervals therefore gives a valid relaxation. A shared-parameter constraint on $z$ cannot be replaced by the same unweighted constraint on occupation masses: choices and visitation probabilities are different objects.

\begin{corollary}[A rational dynamic cost certificate]\label{cor:flow-dual86}
Collect the nonnegative variables in $q$ and write the relaxation as $Aq=b$, $Gq\le h$. Let $w^+$ put the weights $B_td_{t,i,a}^+$ on occupation variables and zero on the other variables. Any rational $y$ and $\lambda\ge0$ satisfying
\[
 A^{\mathsf T}y+G^{\mathsf T}\lambda\ge w^+
\]
certify the upper cost difference $b^{\mathsf T}y+h^{\mathsf T}\lambda$. Applying the same construction to $-B_td_{t,i,a}^-$ certifies the negative lower difference. A numerical solver proposes multipliers; exact checking of every column inequality determines validity.
\end{corollary}

This extension converts the first-date assignment calculation into a dynamic flow calculation. It can improve both the independent-cell bound and the bound obtained by discarding global parameter restrictions. It is sharp for the stated relaxation when strong linear-programming duality applies, not necessarily for the set of policies attainable by a neural network. The distinction is important when fractional indicators or interval transition bounds enlarge that set. Refining cells, parameter sets or transition enclosures reduces this enlargement whenever the old valid restrictions are retained.

\subsection{Recursive preferences and the retained economic extensions}
An expectation-specific occupation identity must not be used as an identity for a nonlinear recursive objective. There is nevertheless a complete-policy transfer for the monotone, cash-invariant evaluators in \eqref{eq:operator49}. For a parameter subset $D\subseteq\Theta_\epsilon$, let $\mathcal A_t^D(x)$ contain every action selected at $(t,x)$ by a policy with parameter in $D$. Define
\begin{align}\label{eq:nonlinear86}
 L_T^D=U_T^D&=g,\\
 L_t^D(x)&=\inf_{a\in\mathcal A_t^D(x)}Q_t(L_{t+1}^D;x,a),\\
 U_t^D(x)&=\sup_{a\in\mathcal A_t^D(x)}Q_t(U_{t+1}^D;x,a).\notag
\end{align}
Measurability and effective near-extrema are assumed as in the original operator. Certified integration and near-extremum errors are charged explicitly.

\begin{proposition}[Nonlinear recursive transfer]\label{prop:nonlinear86}
For every $\theta\in D$, the complete recursive cost obeys
$L_t^D\le V_t^{\pi^\theta}\le U_t^D$ at every date. For any finite cover $\Theta_\epsilon=\bigcup_rD_r$, the bounds $\min_rL_0^{D_r}$ and $\max_rU_0^{D_r}$ enclose all initial recursive costs. Refining the parameter cover and retaining valid action enclosures cannot worsen these exact bounds. If numerical lower and upper recursions incur local errors at most $\eta_t$, expanding either endpoint by $\sum_{s\ge t}(B_s/B_t)\eta_s$ is sufficient.
\end{proposition}

The proof uses monotonicity for the exact bounds and nonexpansiveness for numerical transport. It is supplied in the Technical Supplement. The parameter cover restricts all dates jointly; it does not fit a new continuation independently at each visited state. The expected-cost flow method and the nonlinear envelope method thus serve different objectives within the same NBO construction.

For diffusions, the original discretization and integration allowances enter $\eta_t$. Endogenous preference variables can be included in the state when the stipulated Markov recursion is valid. For temporal selves, the construction applies separately to each self's declared recursive evaluation. In games it applies to each unilateral-deviation problem with opponents fixed. These statements preserve the separate equilibrium and consistency requirements in the retained application results; they do not turn a unilateral cost bound into a proof of equilibrium existence. The recursive-preference, temporal-self and game developments remain part of the paper's mathematical program, rather than being replaced by the two-date investment experiment.
